\subsection{Proof of the $d$-dimensional kernel}
For $1\le d<n$, the $d$th forward difference has the explicit binomial stencil
\begin{equation}\label{eq:stencil2}
 (D_d\tau)_t=\sum_{k=0}^{d}(-1)^{d-k}\binom{d}{k}\tau_{t+k},\quad
 t=0,\ldots,n-d-1.
\end{equation}
The first $n-d$ columns of $D_d$ form an upper triangular
$(n-d)\times(n-d)$ submatrix with nonzero diagonal $(-1)^d$.
Hence $\mathrm{rank}(D_d)=n-d$. As $Q_d=D_d^\top D_d$,
$x^\top Q_dx=\|D_dx\|^2$, and therefore
\begin{equation}\label{eq:kernel-proof}
\ker Q_d=\ker D_d,\qquad
\dim\ker Q_d=n-\mathrm{rank}(D_d)=d.
\end{equation}
The vectors $1,t,\ldots,t^{d-1}$ have vanishing $d$th forward differences
and are linearly independent on $n>d$ nodes. Thus they give a basis
of the kernel, proving that exactly $d$ eigenvalues of $Q_d$ vanish
in exact arithmetic. Under $\mu=0$, the limit $\lambda\to\infty$
is the least-squares projection onto polynomials of degree $\le d-1$.
For instance $d=1$ preserves a constant, $d=2$ a line, and
$d=3$ a quadratic. When $\mu\ne0$, the limiting set is instead
the affine space $\{\tau:D_d\tau=\mu\mathbf1\}$, which has the same
$d$ homogeneous directions but includes a degree-$d$ particular
solution. In the implementation $p_t=\mu\binom{t}{d}$.

\subsection{Interpretation of higher orders and the spline connection}
For a sufficiently smooth function $f$ observed on an equally spaced
mesh of width $h$, Taylor expansion yields
$\Delta_h^df(x_t)=h^df^{(d)}(x_t)+O(h^{d+1})$. Consequently,
under appropriate mesh-refinement assumptions,
\begin{equation}\label{eq:spline-bridge}
h^{1-2d}\sum_t(\Delta_h^df(x_t))^2
\longrightarrow\int |f^{(d)}(x)|^2\,dx.
\end{equation}
This is a mathematical bridge to roughness penalties in smoothing
splines. A Whittaker-type trend smoother penalizes differences
of fitted ordinates directly, while a $P$-spline instead penalizes
differences of coefficients in a chosen spline basis; these are related
ideas but not identical estimators. For a continuous penalty weight
$\lambda_{\rm cont}$, the corresponding discrete coefficient is scaled
approximately as $\lambda_{\rm disc}=\lambda_{\rm cont}h^{1-2d}$.

The order $d$ is an integer representing the order of the derivative
being approximated; it is not a continuous interpolation between
derivative orders. Comparing different $d$ values is a change of
regularization model, so equal numerical $\lambda$ does not
imply equal smoothing. It is useful to examine orders
$d=1,\ldots,5$ under a common normalized trace-based index;
however, arbitrary orders and polynomial nullspaces have long been
known \cite{weinert2007wh,biessy2026wh}. Larger $d$ can be
ill-conditioned because the binomial stencil has growing coefficients
and the positive eigenvalues may span many orders of magnitude.
Neither higher $d$ nor higher polynomial degree guarantees a better
trend or numerical stability.
